\documentclass[12pt]{article}
\usepackage[utf8]{inputenc}
\usepackage[T1]{fontenc}
\usepackage{hyperref}
\title{Burma Pyu Language Resources for K-12 Curriculum Development}
\author{Compiled by Claude Sonnet 5.5 (claude-sonnet-5-5) for Andr\'as Kornai}
\date{}
\begin{document}
\maketitle

\section{Introduction}
Pyu (ISO 639-3 pyx; Glottocode burm1262) is an extinct Sino-Tibetan language of Upper Burma, attested in inscriptions from roughly the 7th to the 12th century. It has no living speakers, so no school curriculum in the language is possible. Only scholarly resources exist.

\section{Primary Educational Resources}
None. No primers, textbooks or school programmes exist for an extinct language.

\section{Linguistic and Reference Materials}
\begin{itemize}
\item Wikipedia article on the Pyu language: \url{https://en.wikipedia.org/wiki/Pyu_language_(Sino-Tibetan)}
\item Marc Miyake, work on Pyu as a distinct branch of Sino-Tibetan (Benjamins catalogue entry): \url{https://www.benjamins.com/catalog/ltba.18013.miy}
\end{itemize}

\section{Religious and Community Resources}
None found.

\section{Digital and Audio Resources}
The search summary mentions the online Corpus of Pyu inscriptions (Griffiths, Miyake and Wheatley, 2018), but no URL for it appeared in the results, so none is given.

\section{Curriculum Development Context}
Pyu could be treated only as heritage content, for example in history or archaeology lessons about the Pyu city-states. This is not a language-medium curriculum.
\end{document}